\chapter{面向雷达资料的来源约束问答应用评价}
\label{chap:radar-application}

\section{研究问题与评价边界}

前述公开实验考察执行反馈恢复轨迹能否改善小模型的知识图谱问答。本章将相同方法链路迁移到雷达资料，分别检查工具接口能否使用、来源记录能否找到，以及模型能否依据原文回答。三个问题需要不同证据：程序执行成功不保证选对记录，选对原文页也不保证正确理解表格中的量值与条件。因此，本章保留中间指标和失败结果，不以一次接口检查替代领域问答评价。

本章沿用公开实验的模型记号：P 为一次生成程序的路线；$A_1,A_2$ 为普通续训路线；$C_1,C_2$ 为恢复轨迹续训路线，下标对应两个续训种子。领域接口适配后仍沿用这些记号，但其权重已发生变化。原始模型在公开留出集上已有的正向结果不能直接赋予新权重，也不能证明 C 在雷达领域优于 A。

本轮领域材料分为两类。合成材料只用于学习字段、条件和工具调用的接口；新来源材料用于一次固定的探索性评价。后者含四份制造商资料及 24 道 AI 编写、AI 核验的问题，尚无人类金标准。它能揭示当前链路的可用性和局限，不能单独建立稳定的领域泛化结论，更不能视为整篇论文已经完成。

\section{来源、知识表示与数据构建}

\subsection{来源范围及可追溯性}

本轮选择四份制造商公开资料，覆盖气象雷达与船用雷达的五个型号。表\ref{tab:domain-sources}列出来源分布，文献\cite{domain-vaisala,domain-leonardo,domain-furuno,domain-jrc}列出文档入口。资料按完整 PDF 页形成 15 个文本块，页块不依赖问题或答案切分；结构化知识库另整理 62 条来源记录，包括 40 条标量、10 条范围、1 条离散选项和 11 条文本记录。每份资料编写六道中文问题，其中八道属于冻结前定义的显式查询条件子集。

\begin{table}[htbp]
\centering\small
\caption{新来源探索材料的范围与历史重叠说明}
\label{tab:domain-sources}
\DomainSourceTable
\end{table}

“新来源”指本轮固定的制造商文档快照及新编问题，不代表此前完全未见这些装备。WRM200、METEOR 735C 所属家族和 WR2120 在历史未使用材料中存在不同程度的出现；JRC 两个型号只是在检查的主要历史文件中未检出，不能证明整个工作区或基础模型预训练中不存在。当前问题没有沿用旧自动问答，但历史问答谱系尚未全部认证。因此，这批材料不能命名为严格未见装备家族测试集。24 道题共享四个来源和家族成分，其相关性也不允许把它们当作 24 个独立来源样本。

PDF 抽取保留可定位的页号、文本哈希和字符区间，部分文档保留版面空格，部分使用阅读顺序文本。抽取仍可能出现连字、空白和表格列对齐问题，原 PDF 的版面是核验依据。结构化记录绑定来源 URL、PDF SHA-256、页号、页块 ID、抽取文本 SHA-256 及字符起止位置；同一型号的不同文档版本不因名称相同而自动合并。公开仓库保存来源索引、协议和聚合结果，不复制资料全文、逐题问答或完整判定内容。

\subsection{量值、选项、条件与来源的联合表示}

雷达参数不是单一的“属性---数字”对应关系。频率、脉宽和脉冲重复频率可以具有不同配置；典型工作距离与技术距离具有不同含义；开机温度与工作温度、天线扫描器与显示单元的条件也不能互换。为保留这些区别，将一条来源记录概括为
\begin{equation}
 r=(e,a,k,v,u,c,t,\pi),
\end{equation}
其中 $e$ 为主体，$a$ 为属性，$k$ 为值类型，$v$ 为值或值集合，$u$ 为单位，$c$ 为原始条件，$t$ 为事件标签，$\pi$ 为来源定位。部件限定保留在主体或条件中，不能在生成回答时丢失。该表示是应用数据的组织方式，本章未将它宣称为已证明优于其他表示的新算法。

标量、连续范围和离散选项分别表示为 $x$、$[l,h]$ 和 $\{x_1,\ldots,x_m\}$；离散选项的最大最小值不构成可任意取值的区间。回答输出合同还支持上界、下界、近似值和文本，保留必要的边界包含性。这些类型是接口能力范围，并不意味着本批 24 道题已经覆盖所有类型。数值和单位可作基础规范化，但规范化不能把不同条件的量值合并，也不能将工程语义相近的词直接当作同一事实。

含引用的回答包含自然语言正文及结构化断言。断言记录主体、属性、值类型、数值或选项、单位、条件、事件与来源页块引用，用于机械检查；自然语言正文另行审阅。空检索只说明当前材料和查询没有返回目标，不能据此断言现实中该参数不存在。引用 ID 存在、JSON 可解析以及断言字段一致，分别对应不同检查，不自动保证正文有据或正确。

\section{直接迁移诊断与共同接口适配}

\subsection{直接迁移未达到可用程度}

在冻结的 12 个已暴露雷达目标上，先以中文和英文问题分别检查原始五个检查点。两种语言是相同目标的成对改写，系统提示、工具定义及知识库说明仍包含双语或中文内容，不能把英文问题视为完整英文交互环境。表\ref{tab:domain-transfer}左侧显示每组仅精确选择 0--3 条记录，大量失败涉及参数签名、条件字段与结束动作。两种语言的总正确数恰好相同，但部分正确目标并不相同；结果没有建立稳定的英文优势。

这一现象说明当前接口不能依赖直接迁移，不能推导出模型永远无法完成雷达问答。继续在这 12 个目标上调整措辞会增加已见材料暴露，因此停止该路线的提示搜索，改用不含真实装备答案的共同接口适配来定位原因。

\subsection{共同合成适配的设置与结果}

合成材料包含 100 个虚构知识实例组，按 60/20/20 划分训练、开发和保留组，每组提供中英文两个视图，对应 120/40/40 条轨迹。实体、知识实例与外层措辞模板按组隔离；任务语义家族、字段、工具序列及部分术语有意共享。数值不承担真实雷达物理含义，不能进入真实装备事实库。保留集检查的是短程接口在新实例和措辞上的使用，而非未见推理结构或新文档泛化。

四个 A/C 检查点使用完全相同的成功轨迹，每条只监督三个助手动作，不监督工具观察。P 使用相同问题实例的一次程序目标。适配从各自已有 LoRA 继续训练，固定五轮、学习率 $5\times10^{-5}$、微批量 2、梯度累积 8，共 40 次优化更新，不按保留集选最佳检查点。两个配对分别沿用种子 20261003 和 20261004，组内核验实际采样顺序相同。每个模型接受 600 次样本呈现；每个 A/C 的实际监督 token 为 48,670，P 为 21,000。因此 A/C 适配预算匹配，P 与 Agent 的监督 token 不匹配。

\begin{table}[htbp]
\centering\small
\caption{共同适配前后记录集合精确选择数；旧题每语言 12 题，合成材料为 20 组的 40 个双语视图}
\label{tab:domain-transfer}
\resizebox{\linewidth}{!}{\DomainTransferTable}
\end{table}

适配后五个模型在合成保留视图上达到 39--40/40，已见雷达题也大幅改善。共同成功轨迹可以修复这一有限接口，是当前结果支持的直接结论。改善同时发生于 A 和 C，不能归因于恢复训练的额外领域优势。近乎满分还受到相同三动作结构和共享术语的限制，不能当作真实事实准确率。旧题的 24 个语言视图已多次暴露，只作为诊断保留。P 的事件选择错误、$A_1$ 的实体拼写和终止失败、$C_2$ 的英文条件遗漏均被保留，未选择有利种子或语言重新报告。

本阶段五次训练及前后推理共计 15 个受监管 GPU 任务，记账约 0.9160 GPU 小时；这一数值不是训练内核时间。适配后的公开数据集保持能力尚未重评，原始公开实验的分数属于原始检查点。

\section{同来源、同生成器的六路线评价}

\subsection{固定证据入口与回答协议}

为区分接口修复和真实资料阅读，本轮对新来源的相同 24 道中文题保留 RAG、P、$A_1$、$C_1$、$A_2$、$C_2$ 全部六组。RAG 直接检索完整原文页块，固定 BM25 参数 $k_1=1.2,b=0.75$，最多返回四块；中文采用汉字一元、二元切分，并使用预先固定的通用中英字段扩展。扩展词不含型号、目标值或逐题来源路由，不进行重排及分数驱动的调整。准备阶段曾观察到检索支持页覆盖 23/24，该观察已在冻结协议中披露，不能称检索设置完全未见本批材料。

P/A/C 使用适配后的检查点选择结构化记录，再将这些记录映射到原文页块，去重、按固定 ID 排序并最多保留四块。所有路线的回答模型都为\emph{同一个未适配的 Qwen2.5-3B-Instruct 基础检查点}，使用同一提示和输出合同，只接收问题、原文页块标题、内容与引用 ID。知识库的结构化值、中文解释和目标记录 ID 不进入回答提示。因而各路线在回答阶段受到一致的原文阅读要求，但知识库构建和选择器训练的额外成本仍然存在，不能声称这是准备条件完全等价的单算法对照。

选择器最多 24 次调用，总生成预算为 2,048 token；逐步 Agent 每轮最多 192 token。回答阶段固定批量 8、上下文上限 8,192 token、最多生成 768 token，均使用贪心解码，输入超限前置拒绝而非静默截断。先完成五组记录选择和六组回答，再对 120 条选择轨迹与 144 条回答输入、解析结果进行 CPU 复核，全部无不一致后访问参考。运行协议、检查点、输入和输出采用哈希绑定；末尾重新核对基础模型权重哈希。本轮仅推理，没有新训练。

\subsection{指标与 AI 审阅口径}

记录精确选择要求最终记录集合与冻结目标集合相同；支持页齐全要求输入包含全部参考支持页。两者不同：选错记录仍可能映射到包含答案的整页，后续回答正确不能反推选择器正确。JSON 合同有效性检查输出格式，严格断言匹配检查规范化后的字段。后者只处理基础单位、大小写和空白等有限变化，不解决同义表达、语义改写或断言拆分方式差异。

自然语言评价将 144 个输出去重为 58 个唯一答案，由两个未见路线标签的 AI 审阅者分别检查正文正确性、所给原文支持、引用支撑和条件保留。正确性判断没有分歧；五个唯一答案在“正文有据”上存在分歧，由已见分组结果的主审依据来源裁决。原判断和分歧均保留。两名 AI 代理不构成不同模型家族共识，也不等同于独立人工核验。

“正文有据”不等于“正文回答正确”：证据不足时，适当的限域拒答可有据但未回答问题。引用支撑也不等于引用 ID 存在。对于题面已明确给出的条件，回答可以承接问题语境；对另需保留的限定、近似或离散选项不能随意省略。表中的条件指标固定使用八道显式查询条件题作为分母，另有九题的参考限定诊断不用于替换该分母；条件保留本身也不保证整句正确。

\subsection{完整结果与成本}

\begin{table}[htbp]
\centering\small
\caption{六路线的新来源结果；除条件列外计数分母均为 24，RAG 无记录选择步骤}
\label{tab:domain-answer}
\resizebox{\linewidth}{!}{\DomainAnswerTable}
\end{table}

\begin{table}[htbp]
\centering\small
\caption{格式、正文支持与引用的分项评价；每列分母为 24}
\label{tab:domain-support}
\resizebox{\linewidth}{!}{\DomainSupportTable}
\end{table}

表\ref{tab:domain-answer}中的正文正确数依次为 13、13、16、15、16、16。两组配对中，$C_1-A_1$ 为 $-1$ 题，$C_2-A_2$ 为 0 题；第二组逐题正确性也相同。该结果没有显示恢复轨迹训练带来额外雷达回答收益。Agent 路线的部分计数高于本轮 RAG 和 P，但样本小、来源相关且准备过程不同，不能据此宣称稳定或因果性的领域优势。

六组严格断言匹配均为 0/24，而自然语言正文有 13--16/24 被判正确。两个指标的任务不同：严格字段匹配失败可能来自术语、断言拆分与条件表示差异；格式有效也可能伴随事实错误。故本章同时保留严格诊断和独立正文审阅，不能把严格匹配的零分写成“模型所有答案都错”。表\ref{tab:domain-support}进一步显示，正文正确和引用支撑存在差别，尤其不能只看 RAG 的正确数而忽略来源引用质量。

\begin{table}[htbp]
\centering\small
\caption{按来源的正文正确数；每个单元格分母为 6，仅为描述性分组}
\label{tab:domain-source-group}
\DomainSourceGroupTable
\end{table}

表\ref{tab:domain-source-group}保留来源分布，避免将相似页面上的多个问题误当独立证据。当前不进行显著性检验，也不选择最有利来源、种子或题型作为主结果。

\begin{table}[htbp]
\centering\small
\caption{完整选择与回答链路成本；GPU 小时来自分配设备的任务存续时间}
\label{tab:domain-cost}
\resizebox{\linewidth}{!}{\DomainCostTable}
\end{table}

六路线共消耗 \DomainTotalTokens{} token，11 个任务合计 \DomainGPUHours{} GPU 小时，低于本轮 3.25 GPU 小时上限。token 包括各轮完整非填充输入及生成 token，重复前缀重复计入，并非 FLOPs。GPU 时间含加载等任务开销，不是纯推理内核利用时间。RAG 的选择 token 为零仅表示 BM25 不调用语言模型，不代表 CPU 检索没有成本；数据整理、审阅及历史训练成本也未计入本轮推理表。

Agent 的回答阶段原文较短，但加上多轮查询轨迹，总 token 均高于本轮 RAG。C 在本轮比配对 A 少用 token，可以作为运行现象记录；它同时没有提高正文正确数，不能写成已建立正确性与效率双重优势。此前接口训练成本也应与本轮推理成本分开报告。

\section{失效分析与论文中的结论边界}

\begin{table}[htbp]
\centering\small
\caption{错误正文的阶段划分；两列之和为各路线的错误数}
\label{tab:domain-failure-stage}
\DomainErrorStageTable
\end{table}

表\ref{tab:domain-failure-stage}显示，$C_2$ 已向回答模型提供全部 24 道题的支持页，正文仍只正确 16 道。四个 Agent 路线共同失败八道题，$C_1$ 另多错一道。这使“继续改善检索即可解决问题”的解释不充分；在当前配置下，原文量值语义、表格配对与描述完整性也是实际瓶颈。不同路线有相同错误不等于证明某种注意力或训练机制，以下仅作事后描述。

\begin{enumerate}
 \item \textbf{量值类型误读。} 将离散脉宽选项写成连续区间，或把典型工作距离和技术距离合并为单一区间，数值出现于原文并不能证明组合后的断言成立。
 \item \textbf{条件与表格配置绑定错误。} 混淆脉宽与 PRF 的配置对应、启动与运行条件，或遗漏功率所需的波束限定。单个数字正确仍可能没有回答指定条件下的问题。
 \item \textbf{主体及描述不完整。} 混用扫描器与显示单元参数，只给器件型号而遗漏题目要求的类别，或者把发射机和调制器描述混合，均影响完整回答。
 \item \textbf{证据与引用链路失败。} 缺少支持页会产生未答、错误补全或引用不足；即使页面在上下文中，引用未指向支持相应陈述的内容，仍不能认定可追溯性合格。
\end{enumerate}

本章证据支持的应用结论是：共同的少量成功轨迹能够使原本不可用的查询接口达到有限任务上的可用程度；来源页和记录的分离评价能够定位“查询正确而阅读错误”及“记录选错但页面仍足够”的差异。公开实验中的恢复轨迹增益与本章中的应用适配属于不同层次，前者不能代替后者。当前没有证据把 C 的领域额外增益、强化学习收益或语义偏好训练成效列为论文结论。

据此，本轮停止追加领域模型训练和本批题上的提示调参，保留全部检查点、错误与零增益结果。接下来的工作是完善数据构建章与应用章衔接，并按来源多样性、条件类型和表格结构预定义后续评价覆盖范围；扩展标准应先于新模型分数确定，不能只补有利样例。若未来讨论同检索结果下原文页、规则模板及条件绑定证据的回答对照，应另建新来源题目和冻结协议；该方向尚未实验，不预先承诺有效。

24 道 AI 题只完成最小探索验证。来源覆盖、原始版面和事实核验、评价者一致性及适配后公开能力保持等工作仍不完整。正式论文还需适配学校模板、补足数据构建说明与可信评价证据。本章可以如实报告负结果和边界，但不能把本轮实验收束等同于应用结论充分或学位论文整体完成。
